\documentclass[11pt]{article}
\usepackage[margin=1in]{geometry}
\usepackage{listings}
\usepackage{xcolor}
\usepackage{enumitem}
\usepackage{amsmath}

\lstset{
  language=Java,
  basicstyle=\ttfamily\small,
  keywordstyle=\bfseries\color{blue!70!black},
  commentstyle=\color{green!50!black},
  stringstyle=\color{red!60!black},
  showstringspaces=false,
  frame=single,
  framesep=6pt,
  xleftmargin=6pt,
  xrightmargin=6pt,
  aboveskip=10pt,
  belowskip=10pt,
}

\pagestyle{empty}

\begin{document}

\begin{center}
  {\Large \textbf{CSCI-UA 102 --- Quiz 2}} \\[6pt]
  {\large Stacks} \\[4pt]
  Name: \underline{\hspace{2.5in}} \quad NetID: \underline{\hspace{1.5in}}
\end{center}

\bigskip

\section*{Asteroid Collision}

An array of integers represents asteroids in a row.
For each asteroid, the \textbf{absolute value} is its size and the \textbf{sign} is its direction:
positive means moving \textbf{right}, negative means moving \textbf{left}.
All asteroids move at the same speed.

When two asteroids moving in \textbf{opposite directions} meet, the smaller one explodes.
If they are the same size, both explode.
Two asteroids moving in the same direction never meet.

\textbf{Write a method} \texttt{int[] asteroidCollision(int[] asteroids)} that returns the state of the asteroids after all collisions, in their original order.

\bigskip

\noindent\textbf{Examples:}
\begin{itemize}[nosep]
  \item \texttt{[5, 10, -5]} $\rightarrow$ \texttt{[5, 10]} \quad (10 and $-5$ collide; $-5$ explodes)
  \item \texttt{[8, -8]} $\rightarrow$ \texttt{[]} \quad (same size, both explode)
  \item \texttt{[-2, -1, 1, 2]} $\rightarrow$ \texttt{[-2, -1, 1, 2]} \quad (no collisions occur)
  \item \texttt{[10, 2, -5]} $\rightarrow$ \texttt{[10]} \quad ($-5$ destroys 2, then loses to 10)
  \item \texttt{[15, 3, 2, 1, -5]} $\rightarrow$ \texttt{[15]} \quad ($-5$ destroys 1, 2, and 3, then loses to 15)
\end{itemize}

\bigskip

\noindent\textbf{You may use the following \texttt{Stack} interface and assume a working implementation exists:}
\begin{lstlisting}
public interface Stack<E> {
    E pop();          // remove and return the top element
    void push(E e);   // add element to the top
    int size();       // number of elements in the stack
}
\end{lstlisting}

\vfill

\noindent\textbf{What is the time complexity of your solution?} \underline{\hspace{2in}}

\newpage

{\color{blue}
\section*{Solution}

\begin{lstlisting}[basicstyle=\ttfamily\small\color{blue},
  keywordstyle=\bfseries\color{blue!70!black},
  commentstyle=\color{blue!50}]
public static int[] asteroidCollision(int[] asteroids) {
    Stack<Integer> s = new LinkedStack<>();
    for (int a : asteroids) {
        boolean alive = true;
        while (alive && a < 0 && s.size() > 0) {
            int top = s.pop();
            if (top < 0) {
                // both going left -- no collision
                s.push(top);
                break;
            }
            // top > 0, a < 0: collision
            if (top > -a) {
                s.push(top);       // top survives
                alive = false;
            } else if (top == -a) {
                alive = false;     // both die
            }
            // else: top < -a, top dies, loop continues
        }
        if (alive) s.push(a);
    }
    int[] result = new int[s.size()];
    for (int i = result.length - 1; i >= 0; i--)
        result[i] = s.pop();
    return result;
}
\end{lstlisting}

\textbf{Time complexity:} $O(n)$ --- each asteroid is pushed and popped at most once.
}

\end{document}
